% !TEX root = ../nc-quadric-classification.tex
\section{多パラメータ skew polynomial 環上の正規二次切断}
\label{sec:multiparameter-quadrics}

ここでは $q_{ij}\in k^\times$、$q_{ji}=q_{ij}^{-1}$、$q_{ii}=1$ とし、
$A=k\langle x_1,x_2,x_3,x_4\rangle/(x_jx_i-q_{ij}x_ix_j)_{i<j}$ を考える。
$Q_r$ ($r\in k^\times$) を、次のパラメータの skew polynomial 環の
$g=x_1x_2+x_3x_4$ による商とする。
\[
 q_{12}=q_{34}=1,\qquad (q_{13},q_{14},q_{23},q_{24})=(r,r^{-1},r^{-1},r).
\]

\begin{theorem}[多パラメータ切断の分類]
\label{thm:multiparameter-quadric-classification}
上の $A$ の正規二次切断から得られる $\HH$ の pointed exact 同値類は、
命題~\ref{prop:nonstandard-rank-three-example} の非 cubic 型一つと、
$Q_r$ の族である。$Q_r$ はすべて cubic 型であり、
\[
 (\qgr Q_r,\pi Q_r)\simeq(\qgr Q_{r'},\pi Q_{r'})
 \quad\Longleftrightarrow\quad
 r'^2=r^2\ \text{または}\ r'^2=r^{-2}
\]
が成り立つ。可換類は $r^2=1$ に対応する。
\end{theorem}

\begin{proof}
正規切断に対してその元の切断環が domain かつ滑らかであることは、
定理~\ref{thm:sign-quadric-normal} の末尾と同じ理由で必要である。
正規化自己同型が対角となる最高単項式の議論も、交換係数が非零ならそのまま使える。
各変数の交換指標を $\rho_i=(q_{ij})_j$ と書く。
$g$ の各単項式は同じ指標 $\chi$ を持ち、
$x_i x_j$ の指標は $\rho_i\rho_j$ である。
等しい指標を持つ変数のブロック内では変数は可換であり、
ブロックごとの基底変更は $A$ の自己同型である。
$\rho\mapsto\chi/\rho$ が固定するブロックには二次形式が現れ、
別のブロックへ移す場合にはブロック対に双線形形式が現れる。
それぞれ対角化または階数標準形にすることで、$g$ は
重複のない変数の平方と、互いに素な変数対の積の和になる。

この形で $g$ に現れない変数があれば商は滑らかでない。
実際、その変数の双対 $X$ は $(A/(g))^!$ で平方零であり、
他の双対生成元と非零スカラーを介して交換する。
このため $(X)$ は非零平方零イデアルである。
$w$ の代表は、$g$ に現れる一つの平方の双対平方、
または一つの積に対応する双対の二項和として取れる。
いずれも同じ変数の積からなるため、$Xw$ と $wX$ はスカラー倍で一致する。
\cite[Lemma 4.11]{MU2022} の正則正規二次元 $w$ を逆にしても消えず、
その局所化の次数零部分にも非零の平方零イデアルを与える。
実際 $Xw\ne0$ なので、$w$ を二文字の和に展開すると
ある双対生成元 $Y$ に対して $XY\ne0$ となり、$XYw^{-1}$ がその部分に入る。
従って $C(B)$ は半単純でなく、\cite[Theorem 5.5]{MU2022} により
商は滑らかでない。

全変数が現れる場合、可能な形は四平方、二平方と一対の積、二対の積である。
四平方なら各成分で $\chi_j=q_{jj}^2=1$ なので全交換係数が符号となり、
定理~\ref{thm:sign-quadric-normal} が適用できる。
二平方の場合に $g=x_1^2+x_2^2+x_3x_4$ とする。
指標の等式から $a=q_{13}$ とすれば
$q_{14}=a^{-1}$、$q_{23}=\pm a$、$q_{24}=q_{23}^{-1}$、
$q_{34}=a^{-2}$、$q_{12}=\pm1$ である。
対角自己同型 $(1,1,a,a^{-1})$ による Zhang twist は
$g$ の直線を保存し、交換係数をすべて符号にする。
従ってこの場合も符号型の分類に帰着する。

二対の積の場合は $g=x_1x_2+x_3x_4$ と取る。
$p=q_{12}$、$q=q_{34}$、$a=q_{13}$ とすると
\[
 q_{14}=p/a,\qquad q_{23}=1/(qa),\qquad q_{24}=qa/p
\]
である。$b^2p=c^2q=1$ を満たす $b,c$ を選び、
対角自己同型 $(b,b^{-1},c,c^{-1})$ で twist すれば、
新しい $p,q$ はともに一になり、$a$ は $r=ba/c$ となる。
二つの積の係数は非零なので変数のスカラー変更で $g$ を正規化できる。
これにより $Q_r$ が得られる。
符号型の混合符号の第三類もこの操作により $r^2=-1$ の $Q_r$ へ移る。

次に全ての $Q_r$ が $\HH$ に属することを示す。
$x_1$ は正則正規であり、それを逆にすると
$x_2=-x_1^{-1}x_3x_4$ を消去できるので、商は domain である。
正則性と消去は定理~\ref{thm:sign-quadric-normal} の議論に従い、
交換係数に上の $r$ を入れて確かめられる。
$r^2=1$ の場合は符号 twist で可換例になる。
$r^2\ne1$ の場合、
\[
 \Phi=\begin{pmatrix}x_1&x_3\\x_4&-r x_2\end{pmatrix},\qquad
 \Psi=\begin{pmatrix}x_2&x_3\\x_4&-r^{-1}x_1\end{pmatrix}
\]
は積が両方向とも $gI$ となる線形行列分解を与える。
可換多項式の行列式は比例しないため、二つの MCM 型は相異なる。
$C(Q_r)$ が $M_2(k)\times M_2(k)$ となることと滑らかさは、
定理~\ref{thm:sign-quadric-normal} と同じ次元の議論で従う。
両方の行列の成分は四次元を張る。
各変数の交換係数の積が一で $g$ は中央なので、
Nakayama 自己同型は $Q_r$ でも恒等である
\cite[Example 5.5 and Lemma 1.5]{RRZ2014}。
従って通常シフトが反 Serre の平方根となる。
飽和性、例外性および豊富性は正則切断の議論で確認できる。
$x_1$ で割った後さらに $x_3$ で割ると二変数の skew polynomial 環が得られ、
その tails の構造対象に無限下降列がある。
従って同じ $x_1$ の冪の議論で Gabriel--Krull 条件も従う。
以上で $\HH$ に入り、定理~\ref{thm:cubic-coverage} から cubic 型となる。

最後に同値を判定する。$r\mapsto-r$ は $g$ の直線を保つ
対角符号 twist で実現でき、$r\mapsto r^{-1}$ は $x_1,x_2$ の交換で実現できる。
必要性には命題~\ref{prop:first-orbit-finite-equivalence} の
添字付き Hom 環の同型を使う。
関係の双線形化が定める点対 $(p,p')\in\PP^3\times\PP^3$ では、
交換関係から $p,p'$ の非零座標の集合が一致する。
さらに三座標が非零なら対応する三角形の交換比の積が一でなければならない。
この積は $r^2$ または $r^{-2}$ なので、$r^2\ne1$ なら非零座標は高々二つである。
$g$ の双線形化は座標対 $12,34$ を排除するため、
各射影の像は座標対 $13,14,23,24$ の四直線からなる正方形である。
$r^2=1$ の可換類では像は滑らかな二次曲面なので、
二つの領域は添字付き線形同型で交わらない。

両方が $r^2\ne1$ の領域にあれば、その同型の各添字の射影線形変換は
四直線の交点、すなわち四座標点を置換する。
点対が各座標点を同じ座標点へ結ぶので、この置換 $\pi$ は添字に依存しない。
従って第 $i$ 添字で $x_a$ を $c_{i,a} x_{\pi(a)}$ へ送ると書ける。
隣接二添字で交換関係を送ったとき、任意の異なる $a,b$ について
$q'_{\pi(a)\pi(b)}/q_{ab}$ は
$c_{i+1,b}/c_{i,b}$ を $c_{i+1,a}/c_{i,a}$ で割った値である。
これは反対辺の対 $12,34$ についても成り立つ。
その二つの単項式は相互に比例するが零ではなく、
各交換関係の像は一方の対の単項式だけを含むからである。
よって三角形を回ると定数倍の比は相殺し、
交換比の三角積が保存される。三角積は $r^2$ と $r^{-2}$ なので
$r'^2=r^2$ または $r'^2=r^{-2}$ が必要となる。
\end{proof}
